% 原始材料：上级目录「成都七中2029届高一上学期半期考试数学卷」图片 1–4（试题）、5–9（答案）。
\documentclass[UTF8,a4paper,12pt]{ctexart}
\usepackage[margin=2cm]{geometry}
\usepackage{amsmath,amssymb,needspace,fancyhdr,tikz}
\pagestyle{fancy}\fancyhf{}\fancyfoot[C]{\thepage}\renewcommand{\headrulewidth}{0pt}
\setlength{\parindent}{0pt}\setlength{\parskip}{4pt}
\newcommand{\subq}[1]{\par（#1）}
\newcommand{\q}[1]{\par\textbf{#1．}}
\newcommand{\opts}[4]{\par\noindent\makebox[.25\linewidth][l]{A．#1}\makebox[.25\linewidth][l]{B．#2}\makebox[.25\linewidth][l]{C．#3}D．#4}
\newcommand{\pick}{（\hspace{1em}）}
\newcommand{\blank}{\underline{\hspace{3cm}}}
\begin{document}
\begin{center}{\LARGE 成都七中2029届高一半期考试数学参考答案}\\[5pt]2026—2027学年度上期\qquad 满分：150分\end{center}
\textbf{一、单项选择题}\quad 1—4：C、A、B、D；\quad 5—8：B、D、A、D。
\q{1}从 $1$ 至 $8$ 中删去 $1,3,5,7,8$，余下 $\{2,4,6\}$。
\q{2}正数平方保序。取 $c=0$ 即排除 B；指数函数 $2^x$ 递增，正数取倒数反序。
\q{3}根号内为 $-3(x-3)^2$，仅在 $x=3$ 时非负。
\q{4}每经过一年乘以 $1+8\%$，故 $y=1000(1+8\%)^x$。
\q{5}$M=[2,+\infty)$，$N=(3,+\infty)$，$N$ 是 $M$ 的真子集。
\q{6}奇性给出 $f(-1)=f(0)=0$，且正半轴亦严格递增。分别与 $-1$、$1$ 处的零值比较，并补上 $0$，得 $(-\infty,-1]\cup[0,1]$。
\q{7}$f(x)=|x|^2+4|x|+1$ 随 $|x|$ 增大而增大。平方得 $4a^2>(a-3)^2$，即 $(a+3)(a-1)>0$。
\q{8}$f(x)=a-(2a+1)/(x+2)$，A 正确。$a=-2$ 时在所给区间递减，左端值为1。$a=2$ 时令 $t=x+2\ge5$，则 $g=2t+10/t-9$；与 $t=5$ 处作差为 $(t-5)(2-2/t)\ge0$。D 中不递减的情形为 $2a+1\ge0$，故选 D。

\textbf{二、多项选择题}\quad 9．AC\qquad 10．ACD\qquad 11．BC
\q{9}$a^x=\sqrt3$，直接代入，A、C 两式左边分别为 $8/3$、2。D 的左边为 $a^{2x}-1+a^{-2x}=7/3$。
\q{10}记 $u=(-3-\sqrt5)/2$，$v=(-3+\sqrt5)/2$。依次在 $u,-1,-1/2,v,0$ 处分段，三函数的中位数从左到右分别是
\[
-x,\quad (x+1)^2,\quad x+1,\quad -x,\quad (x+1)^2,\quad x+1.
\]
各分段端点处相邻式相等。于是 $f(-1)=0$，$f(x)=1$ 仅有 $x=-2,0$。在 $(-1/2,v)$ 上 $f=-x$ 递减；在 $[u,-1]$ 上，D 式右边减左边等于 $(x_1-x_2)^2/4\ge0$。
\q{11}令 $p=xy\in(0,t^2/4]$，则
\[
M=\frac{t^2+2-2p}{t^2+1-2p+p^2},\qquad N=p+\frac{1+t^2}{p}-2.
\]
$N$ 连续且在 $p\to0^+$ 时趋于无穷，故有最小值、无最大值。$t=1$ 时 $M-3/2=p(2-3p)/(2(2-2p+p^2))>0$，下确界 $3/2$ 不能取得。$t=2,4$ 时将 $M$ 连续延拓到 $p=0$，有
\[
M(p)-M(0)=\frac{p[2-(t^2+2)p]}{(t^2+1)(t^2+1-2p+p^2)}.
\]
该差在靠近零时为正，在 $p=t^2/4$ 时为负，故闭区间上的最大、最小值均在 $p>0$ 处取得。
\newpage
\textbf{三、填空题}
\q{12}$\forall a\in\mathbb R$，方程 $x^2-ax+1=0$ 有实根。
\q{13}$f(x)=\dfrac2{x-1}\ (2\le x\le6)$。两点代入得 $2b+2=2(2+a)$、$5(6b+2)=2(6+a)$，故 $a=-1,b=0$。原卷规定漏写定义域扣2分。
\q{14}$\left[1,\frac32\right)\cup[2,2\sqrt2)$。当 $x\to0^+$ 时 $f(x)\to0$，故只需在 $(0,2)$ 上处处 $f(x)<2$ 且有一点 $f(x)\ge1$。若 $a\le0$，取 $x=1.5$ 即违反上界。若 $0<a<2$，左段最大值为 $a^2/4<1$，右段递增趋于 $4-2a$，端点不取，故 $1<4-2a\le2$。若 $2\le a<4$，顶点在区间内，须 $1\le a^2/4<2$。若 $a\ge4$，$f(1)=a-1\ge3$。合并即得。
\par\textbf{四、解答题}\quad 括号内为原卷累计评分点。
\q{15}（1）$A=[2,+\infty)$（2分）。当 $a=4$ 时 $B=(-2,2)$（4分），故 $A\cup B=(-2,+\infty)$（6分）。
\subq{2}若 $a\le0$，$B=\varnothing$，成立（8分）。若 $a>0$，$B=(-\sqrt a,\sqrt a)$（9分），无交点等价于 $\sqrt a\le2$，即 $0<a\le4$（12分）。合并得 $a\in(-\infty,4]$（13分）。
\q{16}（1）设两次所得黄金分别为 $m_1,m_2$ 克（1分）。由 $5a=m_1b$ 得 $m_1=5a/b$（2分）；由 $m_2a=5b$ 得 $m_2=5b/a$（3分）。故 $M=5(a/b+b/a)$（4分）。
\subq{2}黄金质量大于10克（6分）。因 $a,b>0$ 且 $a\ne b$，有 $a/b+b/a>2$（8分），故 $M>10$（9分）。
\subq{3}由 $M\le12.5$ 得 $a/b+b/a\le5/2$（10分）。令 $t=a/b>0$，$t\ne1$（11分），则 $(2t-1)(t-2)\le0$，即 $1/2\le t\le2$（14分）。$t=2$ 可取，最大值为2（15分）。
\q{17}（1）$a^x\le a^{-x}$ 当且仅当 $x\ge0$（2分）；当 $x<0$ 时取 $a^{-x}$（4分）。故 $m(x)=a^{|x|}$（5分）。
\subq{2}（i）$m(x)$ 在 $|x|\to\infty$ 时趋于0（7分），故 $b=2$（9分，直接写对亦得本问4分）。
\par （ii）$F(x)=a^{|x|}+2$（10分），不等式化为 $a^{|k|}>a^{|1-k|}$（11分）。因底数小于1，得 $|k|<|1-k|$（12分）。平方得 $k^2<(1-k)^2$，即 $k<1/2$（14分）。解集为 $(-\infty,1/2)$（15分）。
\newpage
\q{18}（1）函数严格递增（1分）。任取 $1<x_1<x_2$，有
\[
f(x_1)-f(x_2)=(x_1-x_2)\left(1+\frac1{(x_1-1)(x_2-1)}\right)
\]
（2分）。第一因子为负，第二因子为正，故 $f(x_1)<f(x_2)$（3分）。依定义得证（4分）。
\subq{2}定义域为 $\mathbb R\setminus\{1-a\}$（5分），关于原点对称给出 $a=1$（6分）。此时函数为 $x-1/x+1-b$，奇性给出 $b=1$（7分）。代回确为奇函数（8分），故 $a+b=2$（9分）。
\subq{3}$g((1,2])=(a+1,a^2+1]$（10分），$f((2,3])=(1,5/2]$（11分）。题意等价于前者包含于后者（13分），即 $a+1\ge1$ 且 $a^2+1\le5/2$（14分）。结合 $a>1$ 得 $1<a\le\sqrt6/2$（16分），故范围为 $(1,\sqrt6/2]$（17分）。
\q{19}（1）幂函数的系数为1，解 $m^2-m-1=1$ 得 $m=2,-1$（1分）。前者给出 $x^{-1}$，图象不关于 $y$ 轴对称（2分）；后者给出 $f(x)=x^2$，合题意（3分）。
\subq{2}$g(x)=x^2-2(a+1)x+4a$ 的对称轴为 $x=a+1$（4分）。若 $a\le-1$，最小值 $4a=-1$ 无合适解（5分）。若 $-1<a<2$，最小值 $-(a-1)^2=-1$，取 $a=0$（7分）。若 $a\ge2$，最小值 $g(3)=3-2a=-1$，得 $a=2$。综上，$a=0$ 或 $2$（9分）。
\subq{3}令 $y=h(x)=x^2-2(a+1)x+4a+2$。由 $h(y)-h(x)=(y-x)(y+x-2a-2)$ 得
\[
h(h(x))-h(x)=(x-1)(x-2a)(x-2)(x-2a-1)
\]
（11分）。记上式为负的解集为 $S$，须对所有 $t\in[1,2]$ 都有 $[t,t+1]\cap S\ne\varnothing$（12分）。

若 $0<a\le1/2$，则 $S$ 在 $[2,3]$ 内无点，取 $t=2$ 即排除。若 $1/2<a<1$，$S=(1,2a)\cup(2,2a+1)$。$t=1$ 时第一个区间与 $[1,2]$ 相交；$1<t\le2$ 时 $[t,t+1]$ 含有紧邻2的右侧点，故也相交（13分）。若 $a=1$，$S=(1,2)\cup(2,3)$，同理成立（14分）。

若 $a>1$，$S=(1,2)\cup(2a,2a+1)$。当 $t<2$ 时第一个区间已满足要求；当 $t=2$ 时须有 $2a<3$（15分）。因此 $a\in(1/2,3/2)$（17分）。
\par\small\textbf{原卷题源说明}\quad 人教A版高中数学必修第一册：第1题 P13例5，第2题 P43第8题，第3题 P55第2(2)题，第4题 P119第2题，第5题 P9第5(2)题，第7题 P101第7题，第9题 P110第7(2)题，第10题 P68例6及P101第8(2)题，第12题 P35第7(1)题，第13题 P81第5题，第14题 P74第13题，第16题 P49第7题，第17(2)题 P120第9题，第18(2)题 P87第13题。《专题与测试卷》：第14题另见P48备选例1，第19(2)题见P47变式2。
\end{document}
